\documentclass[12pt]{article}
\usepackage[T1]{fontenc}
\usepackage[utf8]{inputenc}
\usepackage{lmodern}
\usepackage{textcomp}
\usepackage{enumitem}
\usepackage{geometry}
\geometry{margin=1in}
\begin{document}
\begin{center}
Answer Key - History - K-12 Level Assessment 17 \
\small (Answers are ordered exactly as in the question paper.)
\end{center}

\section*{Multiple Choice}
\begin{enumerate}[label=\arabic*.]
\item \textbf{Answer:} C
\\ \textit{Options:} A) a tundra environment.; B) a desert environment.; C) a forest and lakeside environment.; D) an equatorial environment.
\\ \textit{Note:} The early Mesolithic Maglemosian culture thrived in a forest and lakeside environment because it relied on the rich natural resources available in these settings, such as abundant game and freshwater food sources, which were essential for their survival and lifestyle.
\item \textbf{Answer:} B
\\ \textit{Options:} A) archaeologists.; B) ethnographers.; C) linguists.; D) paleoanthropologists.
\\ \textit{Note:} Ethnographers are researchers who engage in participant observation within specific communities to understand their cultures, behaviors, and social practices.
\item \textbf{Answer:} D
\\ \textit{Options:} A) beans, millet, and maize.; B) wheat, rice, and peas.; C) sorghum, emmer, and legumes.; D) millet, cabbage, and rice.
\\ \textit{Note:} The excavation of Yang-shao sites in China reveals evidence of early agricultural practices, specifically the cultivation of millet, cabbage, and rice as key domesticated crops that were foundational to ancient Chinese farming and diet.
\item \textbf{Answer:} B
\\ \textit{Options:} A) microcephaly; B) neoteny; C) prognathism; D) supraorbital cortex
\\ \textit{Note:} Neoteny refers to the retention of juvenile features in the adult form, which in hominids is associated with longer developmental periods that may have allowed for increased brain growth and cognitive capabilities.
\item \textbf{Answer:} A
\\ \textit{Options:} A) Jericho in Israel; B) Caral in Peru; C) Stonehenge in England; D) Olmec along the Gulf Coast
\\ \textit{Note:} Jericho, one of the oldest inhabited cities in the world, is recognized for its early development of agricultural practices and social structures, making it a prime example of a complex society emerging around 10,000 BCE.
\end{enumerate}

\section*{True / False}
\begin{enumerate}[label=\arabic*.]
\setcounter{enumi}{5}
\item \textbf{Answer:} True
\\ \textit{Options:} A) True; B) False
\\ \textit{Note:} The recovery of over 6,500 bones from a single hominid species indicates a large sample size, which allows for a better understanding of the species' anatomical diversity, health, and variation within the population.
\item \textbf{Answer:} True
\\ \textit{Options:} A) True; B) False
\\ \textit{Note:} The Cordilleran ice mass was indeed a significant Pleistocene glacier because it covered large areas of the western United States and Canada, shaping the landscape and influencing climate during the last Ice Age.
\item \textbf{Answer:} True
\\ \textit{Options:} A) True; B) False
\\ \textit{Note:} The statement is true because scientific measurements and studies have consistently shown that mean global sea levels have risen approximately 20 centimeters since the late 1800 s due to factors such as thermal expansion of seawater and melting ice sheets, and projections indicate that this trend will continue due to ongoing climate change.
\item \textbf{Answer:} False
\\ \textit{Options:} A) True; B) False
\\ \textit{Note:} A landfill represents a secondary context because it involves the accumulation of human refuse that has been discarded and altered from its original use, rather than preserving items in their primary, unaltered state.
\item \textbf{Answer:} False
\\ \textit{Options:} A) True; B) False
\\ \textit{Note:} The statement is false because the Inca Empire actually reached its height in the 15 th century, around AD 1500, not around AD 200.
\end{enumerate}

\section*{Long Form Response}
\begin{enumerate}[label=\arabic*.]
\setcounter{enumi}{10}
\item \textbf{Answer:} Evidence supporting the idea that Neandertals were hunters includes the discovery of animal bones with stone tool cut marks and a notable absence of animal teeth marks. These cut marks indicate that Neandertals used tools to process the meat from animals they hunted, suggesting they actively participated in hunting activities. The lack of teeth marks further implies that the animals were not scavenged but rather killed and butchered by Neandertals. This evidence collectively highlights the hunting capabilities of Neandertals and their role as active predators in their environment.
\\ \textit{Note:} The presence of stone tool cut marks on animal bones indicates that Neandertals engaged in active hunting and butchering, while the absence of teeth marks suggests they did not scavenge but rather killed the animals themselves.
\item \textbf{Answer:} The Mogollon people decided to leave their mountain homeland primarily due to environmental factors, with drought being a significant influence. Drought led to a scarcity of water and reduced agricultural yields, making it difficult for the Mogollon to sustain their way of life, which heavily relied on farming and foraging. As their crops failed and food sources dwindled, the community faced challenges in supporting its population. Consequently, these harsh environmental conditions pushed the Mogollon to seek new areas with more reliable resources, ultimately leading to their migration away from the mountains they once called home.
\\ \textit{Note:} correct because it identifies drought as a critical environmental factor that led to water scarcity and diminished agricultural production, ultimately forcing the Mogollon people to abandon their mountain homeland in search of more sustainable living conditions.
\end{enumerate}

\vfill
\noindent\textit{End of Answer Key}
\end{document}